\input{../tex-inputs/latex-leseansicht-vorspann}

\section[Arthur Schnitzler an Gustav Schwarzkopf, 28. 7. 1922]{L04530 Arthur Schnitzler an Gustav Schwarzkopf, 28. 7. 1922}
\nopagebreak\mylabel{L04530v}
\rehead{ }\normalsize\beginnumbering\briefempfaengerindex{Schwarzkopf, Gustav@\textsc{Schwarzkopf, Gustav}!zzzSchnitzler, Arthur@\emph{von Arthur Schnitzler}!1922-07-282@{28. 7. 1922}|(be}
\toendnotes[C]{\smallbreak\pagebreak[2]}
\correspDesc{Versand  durch Arthur Schnitzler am 28. 7. 1922 in Feldafing
\newline{}Übermittlung  am 28. 7. 1922 in Feldafing
\newline{}Erhalt  durch Gustav Schwarzkopf im Zeitraum [29. 7. 1922 – 2. 8. 1922?] in Wien}\toendnotes[C]{\smallbreak}
\Standort{DLA, A:Schnitzler, HS.1985.1.1897.}
\physDesc{Bildpostkarte, 409 Zeichen
\newline{}Handschrift: Bleistift, lateinische Kurrent
\newline{}Versand: Stempel: »\nobreak{}\oindex{Feldafing@\textbf{Feldafing}|pwk}Feldafing, 28. Jul 22, 3–4 Nm\nobreak{}«.  }\toendnotes[C]{\smallbreak}\pstart{}{\pb}Herrn Gustav Schwarzkopf\pend{}\pstart{}Wien I\oindex{I., Innere Stadt@\textbf{I., Innere Stadt}|pw}\pend{}\pstart{}Tiefer Graben 17\oindex{Wien@\textbf{Wien}!I., Innere Stadt@\textbf{I., Innere Stadt}!Tiefer Graben 17@\textbf{Tiefer Graben 17}|pw}\pend{}{\bigskip}
\pstart
           \noindent{}\raggedleft{}{\pb}\textcolor{gray}{\textbf{Feldafing am Starnberger See}}\oindex{Feldafing@\textbf{Feldafing}|pw}\pend
           \vspace{1em}
\pstart
           {\pb}28. 7. 22\pend
           \vspace{0.5em}
\pstart
           Herzliche Grüße! trotz Regen und
                    verhängten Himmels ist es in diesem
                    vortrefflichen Hotel\oindex{Hotel Kaiserin Elisabeth [Feldafing]@\textbf{Hotel Kaiserin Elisabeth [Feldafing]}|pwv} (Feldafing, Kaiserin Elisabeth\oindex{Hotel Kaiserin Elisabeth [Feldafing]@\textbf{Hotel Kaiserin Elisabeth [Feldafing]}|pw}) durchaus erträglich; –
                    das Gebirge lockt augenblicklich \textcolor{gray}{wieder}
                    sehr; we{\geminationn} das Wetter sich klärt,
                    will ich mit Heini\pwindex{Schnitzler, Heinrich 9.\,8.\,1902 Hinterbrühl – 12.\,7.\,1982 Wien@\textsc{Schnitzler, Heinrich} (9.\,8.\,1902 Hinterbrühl – 12.\,7.\,1982 Wien)|pw} einen Höhenflug
                    unternehmen, eh er nach Salzburg\oindex{Salzburg@\textbf{Salzburg}|pw}
                    fährt, u mir was über 1000 suchen, \label{T_L04530-1v}\edtext{Anfang August voraussichtlich \label{K_L04530-1v}\edtext{Berchtesgaden\oindex{Berchtesgaden@\textbf{Berchtesgaden}|pw}}{\lemma{\textnormal{\emph{Berchtesgaden}}}\Cendnote{\textnormal{Siehe A. S.: \emph{Wiener Schnitzler}, 10. 8. 1922.}}}\label{K_L04530-1}.}{\lemma{\textnormal{\emph{Anfang … Berchtesgaden.}}}\Cendnote{\textnormal{Der letzte Halbsatz ist entlang des rechten Randes geschrieben.}}}\label{T_L04530-1}\pend
           \selectlanguage{ngerman}\endnumbering\briefempfaengerindex{Schwarzkopf, Gustav@\textsc{Schwarzkopf, Gustav}!zzzSchnitzler, Arthur@\emph{von Arthur Schnitzler}!1922-07-282@{28. 7. 1922}|)be}\mylabel{L04530h}
\begin{anhang}
\end{anhang}\newcommand{\dateiname}{L04530}\newcommand{\titel}{Arthur Schnitzler an Gustav Schwarzkopf, 28. 7. 1922}\newcommand{\editorInnen}{Herausgegeben von Jahnke, SelmaMüller, Martin Anton}\input{../tex-inputs/latex-leseansicht-abspann}
